\chapter{Unsere ostindische Mission}
\quellseite{5}

\section{Nachrichten aus Krishnagiri}

Einem im vorigen Herbst geschriebenen Briefe Missionar Naumanns, der mit Missionar Käther zusammen in Krishnagiri arbeitet, entnehmen wir Folgendes:

Bengadaramen ist ein alter Mann aus dem nahen Pariadorfe. Früher war er einmal Gerber gewesen, jetzt aber verdient er sich seinen Lebensunterhalt mit Strickdrehen. Da seine Frauen — er hatte deren zwei — gestorben sind, er auch sonst keine Verwandten in der Nähe hat, so steht er ganz allein. Schon im letzten Jahre hatte er sich zum Taufunterricht gemeldet. Bisher war er zusammen mit zwei christlichen Lehrersleuten, die aber reformirt waren und sich unserer Kirche anschließen wollten, unterrichtet worden. Lazarus — so hieß der Lehrer — war auch schon alt, arbeitete aber treu und fleißig in der Pariahschule. Ehe er noch mit seiner Frau, die Pacial heißt, aufgenommen werden konnte, nahm ihn Gott schnell und plötzlich zu sich. Er starb, nachdem er nur kurze Zeit an der Wassersucht gelitten hatte, am Herzschlag. Noch bis zum letzten Tage war er in seiner Schule thätig gewesen. Im Missionsgarten, wo auch Missionar Käthers Kinder ruhen, ist er neben seinem Neffen, dem Söhnlein des Lehrers Samuel, christlich zur Ruhe gebettet worden. Er ist nun schon bei Jesu, und sein Leib ruht in seiner Kammer, bis einst der Herr auch ihm zurufen wird: „Lazare, komm heraus!“

Pacial, die Wittwe des Verstorbenen, lernte nun mit dem alten Heiden Bengadaramen weiter, bis sie am Sonntag, den 16. August, bei Gelegenheit unserer Abendmahlsfeier aufgenommen wurde. Zuvor war sie mit Bengadaramen im Wochen-Gottesdienst geprüft worden. Zur besseren Befestigung in der Lehre lernte sie nun noch mit dem alten Bengadaramen zusammen weiter.

Bengadaramen wurde Sonntag, den 30. August, in der Kirche getauft. Er hatte selbst gebeten, daß man ihn nun doch taufen möchte. Damit die heidnischen Schulkinder und die anwesenden Heiden auch aus seinem Munde erführen, um was es sich handelte, wurden an ihn noch einige Fragen gerichtet. Durch seine Antworten legte er ein schönes Zeugniß ab und bewies zugleich, daß er trotz seines Alters doch wohl unterrichtet ist. Lesen kann er ebensowenig wie Pacial, darum hatte er alles durch Vorsagen gelernt. Nach dieser Prüfung empfing er die heilige Taufe und dabei den Namen Simeon.

Gleichsam in der elften Stunde hat der Herr diesen Heiden in seinen Weinberg gedungen. Er lasse ihn seinen Heiland, den er im Glauben erfaßt hat, immer fester fassen. Er lasse ihn beständig bleiben, damit er, wenn sein Stündlein kommt, im Frieden abscheiden und seinen Heiland von Angesicht zu Angesicht schauen kann in ewiger Freude und Seligkeit.

Das ist die erste Heidentaufe in Krishnagiri.

Die Arbeit unserer Missionare in Ostindien ist nicht vergeblich. Der Herr ist mit seinem Wort und wirkt durch dasselbe an der Heiden Herzen. Aber die Frucht des ausgestreuten Samens ist nicht in solchem Maße zu sehen, daß man in Missionsberichten fort und fort hinweisen könnte auf große Schaaren solcher, „die zum Christenthum bekehrt worden wären“. Die Missionsarbeit ist einer Aehrenlese gleich: bald wird hier, bald da
\quellseite{6}
eine Aehre nach vieler Arbeit gesammelt. Das gilt besonders von unserer Heidenmission in Ostindien. Gott stärke unsere Brüder dort zu geduldiger und treuer Arbeit. Gott helfe uns, daß wir wissen, daß sie unsere Boten sind, und daß wir für sie beten. Gott sei ihnen aber auch freundlich und erfreue ihr Herz und richte ihren Muth auf dadurch, daß sie dort Aehren finden und heimbringen können.

Die Gemeinde in Krishnagiri zählt nun zwölf eingeborene Christen: die beiden Lehrersleute mit je drei Kindern, Packial und Simeon. Ein Lehrer gehört noch zu der reformirten Kirche, und ein anderer ist ein junger Muhammedaner, der selbst früher Schüler in der Missionsschule gewesen ist.

Seit einigen Wochen hat sich die Pest in Krishnagiri gezeigt. Etwa zehn Fälle sind bisher in der Stadt selbst vorgekommen. Die Obrigkeit hat durchgreifende Mittel zur Ausrottung dieser Seuche in Anwendung gebracht. Man hat nämlich die ganze Stadt räumen lassen, um alles gründlich rein zu machen. Die Leute müssen außerhalb der Stadt in Hütten wohnen, wozu die Regierung den Aermsten die Matten zur Bedachung gibt. Aller Verkehr und Handel ist dadurch unterbrochen. Natürlich leidet unser Schulunterricht dadurch.

Gott behüte unsere Missionare mit ihren Familien und unsere Christen. Er wende alles zum Besten und zur Förderung seines Reiches unter den Heiden. C. M. Z.
